\uuid{dH0A}
\exo7id{6903}
\titre{exo7 6903}
\auteur{ruette}
\organisation{exo7}
\datecreate{2013-01-24}
\isIndication{false}
\isCorrection{false}
\chapitre{Probabilité discrète}
\sousChapitre{Variable aléatoire discrète}
\cpge{mpsi-mp2i.probabilites-univers-fini.d, pcsi.probabilites-variables-aleatoires.d, ptsi.probabilites-variables-aleatoires.d, mpsi-mp2i.esperance-variance.a, pcsi.esperance-variance.a, ptsi.esperance-variance.a, bcpst1.variable-aleatoire-univers-fini.d}
\module{Probabilité et statistique}
\niveau{L2}
\difficulte{}

\contenu{
\question{
Une usine fabrique des transistors. Chaque transistor a une probabilité de 3\%
d'être défectueux. Quelle est la loi du nombre de transistors défectueux dans 
un lot de 100 transistors ? Que vaut son espérance et que représente-t-elle ?
}
}
